\section{Teoria do Git}

\begin{frame}{O que são commits?}
  \centering
  \only<1>{\includegraphics[width=1\linewidth]{git-commit1.png}}
  \only<2>{\includegraphics[width=1\linewidth]{git-commit2.png}}

  \note{
    Um commit é um registro de uma alteração no histórico do projeto.
    Ele representa um estado do projeto naquele momento e guarda informações como quem fez a alteração, quando ela foi feita e uma mensagem descrevendo o que foi alterado.
    Todo commit também possui um identificador único.
    Esse identificador permite referenciar exatamente aquele commit no histórico do projeto.

    Tá vendo? Só isso. Bem simples.

    E aí, conforme criamos commits e fazendo branches e merges, essas relações começam a formar uma estrutura de grafo.

    Mais especificamente, o histórico de commits do Git forma um grafo acíclico direcionado: cada commit aponta para um ou mais commits anteriores.

    Quem achou que não ia utilizar estrutura de dados para nada?
  }
\end{frame}

\begin{frame}{O que são branches?}
  \centering
  \only<1>{\includegraphics[width=1\linewidth]{git-branches.png}}
  \only<2>{\includegraphics[width=1\linewidth]{git-merge.png}}

  \note{
    Uma branch é uma linha independente de desenvolvimento no histórico do Git.

    Imagine que estamos trabalhando na versão principal do projeto e precisamos desenvolver uma funcionalidade nova.
    Em vez de alterar diretamente a versão principal, criamos uma branch e trabalhamos nela separadamente.

    Enquanto isso, a versão principal pode continuar recebendo outras alterações.
    Quando terminamos a funcionalidade, podemos juntar ambas as linhas novamente mediante um merge.

    O merge registra essa integração criando um novo commit que possui como ancestrais as duas linhas unidas.

    Curiosidade: uma branch é apenas um ponteiro para um commit.
    Conforme novos commits são criados, esse ponteiro avança para acompanhar o histórico daquela linha de desenvolvimento.
  }
\end{frame}

\begin{frame}{O que são tags?}
  \centering
  \includegraphics[width=1\linewidth]{git-tags.png}

  \note{
    Uma tag é um nome dado a um ponto específico do histórico.
    Ela é utilizada principalmente para marcar versões importantes do projeto, como uma versão de lançamento. Diferentemente de uma branch, uma tag não representa uma linha de desenvolvimento.
  }
\end{frame}

\begin{frame}{O fluxo de trabalho do Git}
  \centering
  \includegraphics[width=1\linewidth]{git-workflow.jpg}
\end{frame}

\note{
  Vamos imaginar que estamos trabalhando em um projeto.

  Primeiro, estamos simplesmente editando os arquivos.
  Esse é o Working Directory: o estado atual dos arquivos no nosso computador.

  Agora fizemos três alterações: corrigimos um erro, adicionamos uma funcionalidade e também alteramos a documentação.
  Mas talvez não queiramos registrar as três coisas juntas no mesmo commit.
  Queremos fazer um commit apenas com a correção do erro.

  É aí que entra a Staging Area.
  Ela funciona como uma área de preparação: colocamos nela exatamente as alterações que queremos incluir no próximo commit.

  Quando fazemos o commit, essas alterações passam a fazer parte do histórico do nosso Local Repository.
  Esse histórico fica no nosso computador e podemos consultá-lo sem precisar de internet.
}

\note{
  Depois, quando quisermos compartilhar esses commits, enviamos nosso histórico para um Remote Repository.
  Esse tal de Remote Repository pode ser o GitHub, o GitLab ou qualquer outra plataforma que você estiver utilizando.

  Mas ele nem precisa estar na internet.
  Pode ser simplesmente outro repositório em outra máquina ou até outra pasta no seu próprio computador.
  O Git não sabe nem se importa se existe um GitHub por trás disso.
}